\documentclass{umk-matematika}
\begin{document}

\begin{center}
{\Large\bfseries Практическая работа 09}
\end{center}
\vspace{0.5cm}

\textbf{Задача 1.} Найдите значение производной в точке $x\_0$: $f(x) = x^2,\ x\_0 = 3$\par\vspace{10pt}
\textbf{Задача 2.} Найдите $f(x\_0)$: $f(x) = x^2,\ x\_0 = 3$\par\vspace{10pt}
\textbf{Задача 3.} Запишите общий вид уравнения касательной к графику функции $f(x)$ в точке $x\_0$.\par\vspace{10pt}
\textbf{Задача 4.} Составьте уравнение касательной: $f(x) = x^2,\ x\_0 = 3$\par\vspace{10pt}
\textbf{Задача 5.} Составьте уравнение касательной: $f(x) = x^2 - 4x,\ x\_0 = 1$\par\vspace{10pt}
\textbf{Задача 6.} Составьте уравнение касательной: $f(x) = x^3,\ x\_0 = -1$\par\vspace{10pt}
\textbf{Задача 7.} Составьте уравнение касательной: $f(x) = 2x^2 - 3x + 1,\ x\_0 = 2$\par\vspace{10pt}
\textbf{Задача 8.} Найдите точку на графике $f(x) = x^2$, в которой касательная параллельна прямой $y = 4x - 1$.\par\vspace{10pt}
\textbf{Задача 9.} Составьте уравнение касательной к графику $f(x) = x^2 + 2x$ в точке с абсциссой $x\_0 = -3$.\par\vspace{10pt}
\textbf{Задача 10.} Уклон дорожного покрытия описывается функцией $f(x) = -0{,}5x^2 + 3x$, где x — расстояние от начала участка (м). Найдите уравнение касательной в точке $x\_0 = 2$ (это уклон в данной точке).\par\vspace{10pt}
\newpage
\section*{Ответы}

\begin{enumerate}
  \item Ответ: $f'(3) = 6$
  \item Ответ: $f(3) = 9$
  \item Ответ: $y = f(x\_0) + f'(x\_0)(x - x\_0)$.
  \item Ответ: $y = 6x - 9$
  \item Ответ: $y = -2x - 1$
  \item Ответ: $y = 3x + 2$
  \item Ответ: $y = 5x - 7$
  \item Ответ: точка $(2;\ 4)$.
  \item Ответ: $y = -4x - 9$
  \item Ответ: $y = x + 2$
\end{enumerate}
\end{document}
